% ---------------------------------------------------------------------------
% Study report · The external forcings: air temperature and solar radiation
%
% Build:  pdflatex proxy_forcing_report.tex  (twice, for the ToC)
% Figures and table bodies are read from ../outputs/ and are produced by the
% paired notebook, which must be run first.
%
% Structural draft. The sections below state the report's intended content and
% retain visible placeholders for information and figures that the analysis has
% not yet produced.
% ---------------------------------------------------------------------------
\documentclass[11pt,a4paper]{article}

\newcommand{\runninghead}{The external forcings at Gubbio}

% ---------------------------------------------------------------------------
% reportstyle.tex — shared preamble for the study reports under studies/.
%
% Kept deliberately inside the package set shipped with TeX Live "basic":
% no tcolorbox, no siunitx, no enumitem, no titlesec. Everything below
% compiles with a stock pdflatex installation.
%
% Usage, from a report directory one level below a study folder:
%     % ---------------------------------------------------------------------------
% reportstyle.tex — shared preamble for the study reports under studies/.
%
% Kept deliberately inside the package set shipped with TeX Live "basic":
% no tcolorbox, no siunitx, no enumitem, no titlesec. Everything below
% compiles with a stock pdflatex installation.
%
% Usage, from a report directory one level below a study folder:
%     % ---------------------------------------------------------------------------
% reportstyle.tex — shared preamble for the study reports under studies/.
%
% Kept deliberately inside the package set shipped with TeX Live "basic":
% no tcolorbox, no siunitx, no enumitem, no titlesec. Everything below
% compiles with a stock pdflatex installation.
%
% Usage, from a report directory one level below a study folder:
%     \input{../../_shared/reportstyle.tex}
%     \graphicspath{{../outputs/}}
% ---------------------------------------------------------------------------

\usepackage[T1]{fontenc}
\usepackage[utf8]{inputenc}
\usepackage{textcomp}
\usepackage{lmodern}
\usepackage{microtype}
\usepackage[margin=2.4cm,headheight=15pt]{geometry}
\usepackage{graphicx}
\usepackage{booktabs}
\usepackage{tabularx}
\usepackage{array}
\usepackage{amsmath}
\usepackage{xcolor}
\usepackage[font=small,labelfont=bf,justification=justified]{caption}
\usepackage{float}
\usepackage{fancyhdr}
\usepackage[hidelinks]{hyperref}

% --- Palette ---------------------------------------------------------------
\definecolor{rulegrey}{HTML}{B9C0C7}
\definecolor{fillgrey}{HTML}{F2F4F6}
\definecolor{failred}{HTML}{A5202B}
\definecolor{passgreen}{HTML}{1B6B3A}
\definecolor{accent}{HTML}{1F4E79}

% --- Semantic shorthands ---------------------------------------------------
\newcommand{\degC}{\textdegree C}
\newcommand{\mdegC}{mdeg/\textdegree C}
\newcommand{\fail}{\textcolor{failred}{\textbf{fail}}}
\newcommand{\pass}{\textcolor{passgreen}{\textbf{pass}}}
\newcommand{\worse}{\textcolor{failred}{worse}}
\newcommand{\better}{\textcolor{passgreen}{better}}

% --- Highlighted panel -----------------------------------------------------
% #1 = heading, #2 = body
\newcommand{\panel}[2]{%
  \begin{center}%
  \setlength{\fboxsep}{9pt}%
  \fcolorbox{rulegrey}{fillgrey}{%
    \begin{minipage}{0.93\textwidth}%
      {\bfseries\color{accent}#1}\par\medskip
      #2
    \end{minipage}}%
  \end{center}}

% --- Numeric column: right aligned, fixed width ----------------------------
\newcolumntype{R}[1]{>{\raggedleft\arraybackslash}p{#1}}
\newcolumntype{L}[1]{>{\raggedright\arraybackslash}p{#1}}
\newcolumntype{Y}{>{\raggedright\arraybackslash}X}

% --- Table look ------------------------------------------------------------
\renewcommand{\arraystretch}{1.15}
\setlength{\tabcolsep}{6pt}

% --- Headers ---------------------------------------------------------------
\pagestyle{fancy}
\fancyhf{}
\fancyhead[L]{\small\itshape\runninghead}
\fancyhead[R]{\small\thepage}
\renewcommand{\headrulewidth}{0.4pt}

% --- Section spacing without titlesec --------------------------------------
\makeatletter
\renewcommand\section{\@startsection{section}{1}{\z@}%
  {-3.2ex \@plus -1ex \@minus -.2ex}{2.0ex \@plus.2ex}%
  {\normalfont\Large\bfseries\color{accent}}}
\renewcommand\subsection{\@startsection{subsection}{2}{\z@}%
  {-2.6ex \@plus -1ex \@minus -.2ex}{1.4ex \@plus.2ex}%
  {\normalfont\large\bfseries}}
\makeatother

% --- Figure defaults -------------------------------------------------------
\setkeys{Gin}{width=\linewidth}
\captionsetup[figure]{skip=6pt}

%     \graphicspath{{../outputs/}}
% ---------------------------------------------------------------------------

\usepackage[T1]{fontenc}
\usepackage[utf8]{inputenc}
\usepackage{textcomp}
\usepackage{lmodern}
\usepackage{microtype}
\usepackage[margin=2.4cm,headheight=15pt]{geometry}
\usepackage{graphicx}
\usepackage{booktabs}
\usepackage{tabularx}
\usepackage{array}
\usepackage{amsmath}
\usepackage{xcolor}
\usepackage[font=small,labelfont=bf,justification=justified]{caption}
\usepackage{float}
\usepackage{fancyhdr}
\usepackage[hidelinks]{hyperref}

% --- Palette ---------------------------------------------------------------
\definecolor{rulegrey}{HTML}{B9C0C7}
\definecolor{fillgrey}{HTML}{F2F4F6}
\definecolor{failred}{HTML}{A5202B}
\definecolor{passgreen}{HTML}{1B6B3A}
\definecolor{accent}{HTML}{1F4E79}

% --- Semantic shorthands ---------------------------------------------------
\newcommand{\degC}{\textdegree C}
\newcommand{\mdegC}{mdeg/\textdegree C}
\newcommand{\fail}{\textcolor{failred}{\textbf{fail}}}
\newcommand{\pass}{\textcolor{passgreen}{\textbf{pass}}}
\newcommand{\worse}{\textcolor{failred}{worse}}
\newcommand{\better}{\textcolor{passgreen}{better}}

% --- Highlighted panel -----------------------------------------------------
% #1 = heading, #2 = body
\newcommand{\panel}[2]{%
  \begin{center}%
  \setlength{\fboxsep}{9pt}%
  \fcolorbox{rulegrey}{fillgrey}{%
    \begin{minipage}{0.93\textwidth}%
      {\bfseries\color{accent}#1}\par\medskip
      #2
    \end{minipage}}%
  \end{center}}

% --- Numeric column: right aligned, fixed width ----------------------------
\newcolumntype{R}[1]{>{\raggedleft\arraybackslash}p{#1}}
\newcolumntype{L}[1]{>{\raggedright\arraybackslash}p{#1}}
\newcolumntype{Y}{>{\raggedright\arraybackslash}X}

% --- Table look ------------------------------------------------------------
\renewcommand{\arraystretch}{1.15}
\setlength{\tabcolsep}{6pt}

% --- Headers ---------------------------------------------------------------
\pagestyle{fancy}
\fancyhf{}
\fancyhead[L]{\small\itshape\runninghead}
\fancyhead[R]{\small\thepage}
\renewcommand{\headrulewidth}{0.4pt}

% --- Section spacing without titlesec --------------------------------------
\makeatletter
\renewcommand\section{\@startsection{section}{1}{\z@}%
  {-3.2ex \@plus -1ex \@minus -.2ex}{2.0ex \@plus.2ex}%
  {\normalfont\Large\bfseries\color{accent}}}
\renewcommand\subsection{\@startsection{subsection}{2}{\z@}%
  {-2.6ex \@plus -1ex \@minus -.2ex}{1.4ex \@plus.2ex}%
  {\normalfont\large\bfseries}}
\makeatother

% --- Figure defaults -------------------------------------------------------
\setkeys{Gin}{width=\linewidth}
\captionsetup[figure]{skip=6pt}

%     \graphicspath{{../outputs/}}
% ---------------------------------------------------------------------------

\usepackage[T1]{fontenc}
\usepackage[utf8]{inputenc}
\usepackage{textcomp}
\usepackage{lmodern}
\usepackage{microtype}
\usepackage[margin=2.4cm,headheight=15pt]{geometry}
\usepackage{graphicx}
\usepackage{booktabs}
\usepackage{tabularx}
\usepackage{array}
\usepackage{amsmath}
\usepackage{xcolor}
\usepackage[font=small,labelfont=bf,justification=justified]{caption}
\usepackage{float}
\usepackage{fancyhdr}
\usepackage[hidelinks]{hyperref}

% --- Palette ---------------------------------------------------------------
\definecolor{rulegrey}{HTML}{B9C0C7}
\definecolor{fillgrey}{HTML}{F2F4F6}
\definecolor{failred}{HTML}{A5202B}
\definecolor{passgreen}{HTML}{1B6B3A}
\definecolor{accent}{HTML}{1F4E79}

% --- Semantic shorthands ---------------------------------------------------
\newcommand{\degC}{\textdegree C}
\newcommand{\mdegC}{mdeg/\textdegree C}
\newcommand{\fail}{\textcolor{failred}{\textbf{fail}}}
\newcommand{\pass}{\textcolor{passgreen}{\textbf{pass}}}
\newcommand{\worse}{\textcolor{failred}{worse}}
\newcommand{\better}{\textcolor{passgreen}{better}}

% --- Highlighted panel -----------------------------------------------------
% #1 = heading, #2 = body
\newcommand{\panel}[2]{%
  \begin{center}%
  \setlength{\fboxsep}{9pt}%
  \fcolorbox{rulegrey}{fillgrey}{%
    \begin{minipage}{0.93\textwidth}%
      {\bfseries\color{accent}#1}\par\medskip
      #2
    \end{minipage}}%
  \end{center}}

% --- Numeric column: right aligned, fixed width ----------------------------
\newcolumntype{R}[1]{>{\raggedleft\arraybackslash}p{#1}}
\newcolumntype{L}[1]{>{\raggedright\arraybackslash}p{#1}}
\newcolumntype{Y}{>{\raggedright\arraybackslash}X}

% --- Table look ------------------------------------------------------------
\renewcommand{\arraystretch}{1.15}
\setlength{\tabcolsep}{6pt}

% --- Headers ---------------------------------------------------------------
\pagestyle{fancy}
\fancyhf{}
\fancyhead[L]{\small\itshape\runninghead}
\fancyhead[R]{\small\thepage}
\renewcommand{\headrulewidth}{0.4pt}

% --- Section spacing without titlesec --------------------------------------
\makeatletter
\renewcommand\section{\@startsection{section}{1}{\z@}%
  {-3.2ex \@plus -1ex \@minus -.2ex}{2.0ex \@plus.2ex}%
  {\normalfont\Large\bfseries\color{accent}}}
\renewcommand\subsection{\@startsection{subsection}{2}{\z@}%
  {-2.6ex \@plus -1ex \@minus -.2ex}{1.4ex \@plus.2ex}%
  {\normalfont\large\bfseries}}
\makeatother

% --- Figure defaults -------------------------------------------------------
\setkeys{Gin}{width=\linewidth}
\captionsetup[figure]{skip=6pt}

\graphicspath{{../outputs/}}

\title{\vspace{-1.2cm}\textbf{The external forcings: air temperature and solar
radiation}\\[2mm] \large Characterisation of ERA5 and the Gubbio meteorological station, and their compatibility with the on-structure measurements}
\author{Study report · \texttt{studies/proxy\_forcing\_characterization/}}
\date{}

\begin{document}
\maketitle
\thispagestyle{fancy}

\begin{abstract}
\noindent
Air temperature and solar radiation are candidate environmental drivers of the monitored wall response, but their on-structure records do not cover the full monitoring period. This report characterises two longer external sources---ERA5 reanalysis retrieved through Oikolab and observations from the Meteosystem station in Gubbio---before comparing them with the final valid on-structure values produced by Study~01. The comparison is restricted to air temperature and solar radiation and asks whether either external source is sufficiently compatible with the local measurements to substitute for them where the local record is absent. Compatibility is assessed without treating the three sources as physically equivalent measurements.
\end{abstract}

\tableofcontents

% ===========================================================================
\section{Introduction}
\label{sec:introduction}

The structural interpretation developed after Study~01 requires environmental
series that extend across the monitoring period. The measurements made at the
wall cannot provide that continuity. The on-structure air-temperature channel
contains substantial gaps, while solar radiation was introduced only with the
current instrument package and includes periods rejected by the quality tests in
Study~01. External data are therefore considered as possible substitutes, but
their greater coverage does not by itself establish that they represent the
thermal forcing at the wall.

This report considers two external sources. ERA5 is a gridded atmospheric
reanalysis and represents an area rather than the monitored surface. The
Meteosystem record is produced by a meteorological station in Gubbio and is a
point observation, but it is not made at the wall and its exposure differs from
that of the on-structure instruments. Each source is first characterised on its
own terms, including its provenance, variables, coverage, complete time series,
and summer and winter mean daily cycles. A later comparison then brings ERA5,
the meteorological station, and the on-structure measurements onto a common
hourly grid.

\begin{table}[H]
\centering
\small
\begin{tabularx}{\textwidth}{L{2.1cm}L{3.0cm}L{2.0cm}L{2.4cm}Y}
\toprule
\textbf{Source} & \textbf{Nature} & \textbf{Native interval} &
\textbf{Documented extent} & \textbf{Role in this study} \\
\midrule
ERA5 through Oikolab & Gridded atmospheric reanalysis, representing a grid-cell
average around the monitoring coordinates & 1 hour &
1 Jan 2018--12 Aug 2026 & Long-coverage proxy for air temperature and solar
radiation \\
Meteosystem Gubbio & Ground-based meteorological station in Gubbio; a point
measurement away from the monitored wall & 30 minutes &
17 Oct 2018--12 Aug 2026 & Local meteorological proxy for air temperature and
solar radiation \\
\bottomrule
\end{tabularx}
\caption{The two external sources. The extents are those recorded in the current
repository data dictionary and will be recomputed from the input files when the
study is executed. Coverage percentages are reported separately for each
variable because a source may exist while an individual channel is missing.}
\label{tab:external_sources}
\end{table}


% ===========================================================================
\section{ERA5 reanalysis}
\label{sec:era5}

\subsection{Nature and coverage}

ERA5 combines numerical weather prediction with assimilated observations to
produce a spatially continuous reconstruction of atmospheric conditions. The
series used here was extracted at the monitoring coordinates
($43.353430^{\circ}$ N, $12.582047^{\circ}$ E) and is described in the
repository as representing a grid cell of approximately nine kilometres. It is
hourly and currently spans 1 January 2018 to 12 August 2026. The repository data
dictionary reports approximately 99.5\,\% coverage for the ERA5 variables used
by this study; the executed analysis will replace that provisional figure with
counts measured directly from the input file.

The gridded value is not expected to reproduce the microclimate of a
sun-exposed instrument enclosure or masonry surface. Its possible value is
instead continuity: if its temperature and radiation retain a stable relation
with the on-structure measurements over their shared window, ERA5 may provide a
consistent forcing proxy outside that window.

\subsection{Retrieval and repository organisation}

The paired notebook and script \texttt{auxiliary/oiko.ipynb} and
\texttt{auxiliary/oiko.py} retrieve the data from the authenticated Oikolab
weather API. The script requests hourly means in six-month chunks, concatenates
the responses, removes duplicate timestamps and provider metadata columns, and
appends only dates later than the current local export. The consolidated file is
written to \texttt{data/raw/proxies/oikolab\_weather.csv}; Study~02 reads that
file without making a new network request. Timestamps are supplied in UTC.

\subsection{Variables}

Table~\ref{tab:era5_variables} records the variables requested by the retrieval
script that are retained for this report and their interpretation. Every
variable in the table is included in the source characterisation. Only air
temperature and surface solar radiation subsequently enter the three-source
comparison.

\begin{table}[H]
\centering
\small
\begin{tabularx}{\textwidth}{L{4.4cm}L{1.7cm}Y}
\toprule
\textbf{Oikolab variable} & \textbf{Unit} & \textbf{Meaning} \\
\midrule
\texttt{temperature} & \degC & Near-surface air temperature; compared with the
on-structure and Meteosystem air-temperature series. \\
\texttt{relative\_humidity} & fraction & Relative humidity. It is multiplied by
100 on loading when expressed as a fraction. \\
\texttt{dewpoint\_temperature} & \degC & Temperature at which the represented
air mass reaches saturation. \\
\texttt{surface\_solar\_radiation} & W/m$^2$ & Downward short-wave solar
radiation at the surface; compared with the other radiation series. \\
\texttt{total\_precipitation} & mm water equivalent & Accumulated liquid and
solid precipitation expressed as water depth. \\
\texttt{wind\_speed} & m/s & Magnitude of the near-surface wind. \\
\texttt{wind\_direction} & [angular convention required] & Direction from
which the wind originates. It will be normalised to degrees and treated as a
circular variable. \\
\bottomrule
\end{tabularx}
\caption{ERA5 variables retained for characterisation. The unit, zero direction,
and rotation convention of wind direction must be confirmed from the API
response metadata before the table is finalised.}
\label{tab:era5_variables}
\end{table}

\subsection{Complete time series}

The complete ERA5 record is shown before any restriction to the on-structure
measurement window. The seven retained variables occupy separate vertically
aligned panels with a shared time axis, each panel named by its own title and
carrying its unit on the axis. They are ordered air temperature, relative
humidity, solar radiation, wind speed, wind direction, dew-point temperature and
total precipitation: the first five are the channels the daily-cycle figures
draw, in the same order, so that a reader moving between the figures of this
report finds each channel where it was left. Gaps remain visible rather than
being connected across, and wind direction is drawn on a circular angular scale.
This view establishes the actual extent, missing intervals, long-term range, and
seasonal envelope of every retained variable.

\begin{figure}[H]
\centering
\includegraphics[width=\textwidth]{PF_F01_era5_series.png}
\caption{Complete ERA5 record, one panel for each of the seven retained
variables on a shared time axis. Wind direction is drawn as its daily circular
mean, since eight years of hourly bearings fill a panel solid.}
\label{fig:era5_series}
\end{figure}

\subsection{Summer and winter mean daily cycles}

Mean daily cycles are computed separately for summer and winter from the
complete valid ERA5 record. Summer is provisionally June--August and winter is
provisionally December--February, matching Study~01. Each mean is supported by
its contributing day count so that a poorly covered hour is not presented as
equally certain to a well-covered one. Scalar variables use ordinary hourly
means. Wind direction uses a circular hourly mean after its angles have been
normalised to one documented convention.

The two seasons are drawn in one figure rather than two, a channel to a row and
a season to a column, and the two panels of a row share one vertical scale. The
seasonal difference in a channel's day is then a difference in the drawn shape,
rather than one a reader must reconstruct from two axes with different limits on
two separate pages. Five of the seven retained channels are drawn: the two
forcings this study is about, the humidity that accompanies them, and the wind
that ventilates the surface. Dew point and precipitation keep their place in the
complete-record figure of Section~\ref{sec:era5} and in the tabulated diurnal
profile, both of which cover every retained channel; what is narrowed here is the
figure, not the characterisation.

\begin{figure}[H]
\centering
\includegraphics[width=\textwidth]{PF_F02_era5_diurnal_seasons.png}
\caption{ERA5 mean daily cycles, by channel and season. Air temperature,
relative humidity, solar radiation, wind speed and wind direction, one to a row,
with summer in the left column and winter in the right. Every contributing day
is drawn in gray behind its mean, and each panel names the days its mean rests on
and the peak-to-trough amplitude of that mean.}
\label{fig:era5_diurnal_seasons}
\end{figure}

% ===========================================================================
\section{The Gubbio meteorological station}
\label{sec:ground_station}

\subsection{Nature and coverage}

The Meteosystem source is a ground-based station in Gubbio. In contrast to the
spatial average represented by ERA5, it records conditions at one location.
That location is closer to the monitored structure than a regional grid is, but
distance alone does not establish representativeness: elevation, surrounding
surfaces, shading, sensor height, and screen exposure may all separate the
station record from the microclimate at the wall.

The downloaded record is natively half-hourly and currently spans 17 October
2018 to 12 August 2026. The repository data dictionary reports 88.6\,\% coverage
for air temperature and 88.7\,\% for solar radiation. These values will be
recomputed by the study because the source is updated month by month and the two
variables need not share the same missing intervals.

\panel{Information required}{Insert the station's coordinates, elevation,
distance and height difference from the monitored wall, sensor models,
meteorological exposure, calibration information, operator, and maintenance or
relocation history.}

\subsection{Retrieval and repository organisation}

The paired notebook and script \texttt{auxiliary/meteosystem\_italy.ipynb} and
\texttt{auxiliary/meteosystem\_italy.py} retrieve one CSV file per calendar
month from the Meteosystem website. Each response is retained in
\texttt{auxiliary/data/raw/proxies/\_cache/gubbio/}, allowing later cleaning and
assembly without another download. The monthly files are concatenated,
deduplicated, and converted from Italian civil time to UTC. The current
preparation script removes all categorical direction fields before writing
\texttt{data/raw/proxies/meteosystem\_gubbio.csv}. For this characterisation it
will be revised to retain \texttt{Dir}, convert its compass sectors to a
documented angular convention, and remove only the unselected fields. Study~02
aggregates scalar half-hourly observations to hourly means and computes wind
direction by a circular hourly mean.

\subsection{Variables}

Table~\ref{tab:ground_variables} describes the fields supplied by the station.
Every field in the table is included in the station characterisation. Air
temperature and solar radiation alone enter the later compatibility study with
ERA5 and the on-structure measurements.

\begin{table}[H]
\centering
\small
\begin{tabularx}{\textwidth}{L{2.4cm}L{1.7cm}Y}
\toprule
\textbf{Raw field} & \textbf{Unit} & \textbf{Meaning} \\
\midrule
\texttt{Temp} & \degC & Air temperature; compared with ERA5 and the
on-structure air-temperature series. \\
\texttt{Umid} & \% & Relative humidity. \\
\texttt{Dew pt} & \degC & Dew-point temperature. \\
\texttt{Vento} & m/s & Mean wind speed. \\
\texttt{Dir} & compass sector & Mean wind direction. Compass sectors will be
converted to angles and treated as circular data. \\
\texttt{Press} & hPa & Atmospheric pressure. \\
\texttt{Pioggia} & mm & Precipitation depth. \\
\texttt{Int.Pio.} & [unit required] & Precipitation intensity. \\
\texttt{Rad.Sol.} & W/m$^2$ & Solar radiation; compared with ERA5 and the
on-structure radiation series. \\
\bottomrule
\end{tabularx}
\caption{Variables in the Meteosystem monthly files. The reporting interval and
unit of precipitation intensity require confirmation from the station
documentation.}
\label{tab:ground_variables}
\end{table}

\subsection{Complete time series}

The complete station record is shown independently of the shorter on-structure
record, on the layout and the channel order used for ERA5 in
Figure~\ref{fig:era5_series}: air temperature, relative humidity, solar
radiation, wind speed, wind direction, dew-point temperature, precipitation
depth and atmospheric pressure, each panel named by its own title and carrying
its unit on the axis. The plot preserves missing intervals and uses the
half-hourly observations after the quality-control decisions that produce the
consolidated source file. Wind direction is shown on a circular angular scale.

Precipitation intensity is the ninth retained channel and is not drawn. It is
the precipitation depth of the panel beside it divided by the reporting
interval, so the two panels would carry one measurement at two scales. It
remains in the channel inventory and in the tabulated diurnal profile, both of
which cover all nine.

The precipitation panel is drawn on a clipped vertical axis. The station's
precipitation depth contains a single reading of $1.70\times10^{9}$\,mm in one
hour, and its intensity channel contains five readings of the same order; the
largest depth that is physically possible is the next value down, at
30.1\,mm. An axis scaled to the defective sample renders the entire eight-year
record as a flat line beneath one spike, so the axis is bounded by the values
inside the channel's documented plausible range and the panel carries a note, in
the accent colour, stating how many samples fall outside it and how far. The
clipping is a display choice: no value is removed from the data, and the
implausible readings are counted and reported in the channel inventory
(Table~\ref{tab:ground_variables} describes the fields; the executed inventory
records one implausible depth and five implausible intensities).

\begin{figure}[H]
\centering
\includegraphics[width=\textwidth]{PF_F04_gs_series.png}
\caption{Complete Meteosystem record, one panel for each of eight retained
variables on a shared time axis, in the channel order used for ERA5.
Precipitation intensity is omitted as a rescaling of the precipitation depth
above it. The precipitation axis is clipped to the channel's plausible range,
with a note giving the sample it excludes. Wind direction is drawn as its daily
circular mean.}
\label{fig:gs_series}
\end{figure}

\subsection{Summer and winter mean daily cycles}

The seasonal profiles are computed from the complete valid station record,
before restricting it to timestamps shared with the on-structure instruments.
The season definitions and display conventions match those used for ERA5, so
differences between the two characterisations reflect the records rather than
different summaries. Scalar variables use ordinary hourly means. Wind direction
uses a circular hourly mean of the converted angles.

The layout is the one used for ERA5 in Figure~\ref{fig:era5_diurnal_seasons} and
draws the same five channels in the same order, so that the two sources can be
read against each other panel by panel. The station's four remaining
channels---dew-point temperature, atmospheric pressure, precipitation depth and
precipitation intensity---keep their place in the complete-record figure of this
section and in the tabulated diurnal profile, both of which cover all nine.

\begin{figure}[H]
\centering
\includegraphics[width=\textwidth]{PF_F05_gs_diurnal_seasons.png}
\caption{Meteosystem mean daily cycles, by channel and season. Air temperature,
relative humidity, solar radiation, wind speed and wind direction, one to a row,
with summer in the left column and winter in the right, on the layout and the
channel order used for ERA5. Every contributing day is drawn in gray behind its
mean, and each panel names the days its mean rests on and the peak-to-trough
amplitude of that mean.}
\label{fig:gs_diurnal_seasons}
\end{figure}

% ===========================================================================
\section{The two external sources against each other}
\label{sec:cross_source}

Both external sources are now on the common hourly grid, and each has been
characterised on its own terms. Before either is measured against the wall, they
are measured against each other, over every quantity they both report.

This comparison answers a question the on-structure comparison cannot. The wall
record is short, and for radiation it is trustworthy over a narrower window
still, so a disagreement found there says little about how the two sources
behave outside it. Against each other they have their entire shared record and
every quantity in it, not only the two forcings. What that establishes is the
kind of source each one is: where a grid-cell average of roughly nine kilometres
departs from a point observation, in which direction, and whether the departure
is the same in summer as in winter. A quantity on which the two agree closely is
one where the later choice between them will not matter much; a quantity on which
they diverge is one where the verdict of
Section~\ref{sec:comparison} has to be earned rather than assumed.

Agreement is reported against the station, so that a positive bias reads as the
reanalysis running high with respect to the local observation. Neither source is
a reference in the sense the on-structure record is: this is a comparison of two
proxies, not the validation of one against a truth. Every pair uses its own
overlap, and no value is interpolated to manufacture one. Each quantity is
reported over the whole shared record and then separately for each season,
because a bias that is stable across the year and a bias that changes sign with
the season are indistinguishable in a statistic pooled over both.

The scores are computed on values inside each quantity's documented plausible
range. This is the one place in the study where the implausible readings of
Section~\ref{sec:ground_station} are set aside rather than merely counted, and
the reason is that a score cannot survive them: a single rainfall reading of
order a billion millimetres dominates every mean square it enters, and the
resulting row would look like a result while being an artefact of one defective
sample. Nothing is deleted. The readings remain in the frame every other step of
the study uses, remain counted in the channel inventory, and remain visible in
the complete-record figure; what is excluded here is their effect on a
correlation and a root mean square, and the number of values excluded is
reported with the table.

Precipitation carries one further exclusion, and it is of a different kind. The
station's heaviest recorded hour, 30.1\,mm, is real weather: it is well inside
the plausible range, roughly twice the next largest hour in eight years, and is
the sole occupant of the upper two thirds of its panel's axis. For this
comparison alone the precipitation channels are bounded at 20\,mm in the hour,
which excludes that single observation, so that the row and the panel describe
the behaviour the two sources ordinarily produce rather than one storm. This is
a decision about what the comparison covers, not the removal of a defect, and it
is stated here because a reader who assumed otherwise would misread the
precipitation row. The hour is retained everywhere else in this report, and the
bound is a parameter of the study: widening it puts the observation back.

Wind direction is differenced on the circle, as everywhere else in this study.
Its bias and correlation are the columns to read on that row: the amplitude ratio
and the phase lag are computed by the same shared machinery as for the scalar
quantities, where for a bearing they do not carry a meaningful peak-to-trough
distance, and the ratio is reported as missing. It is missing for precipitation
too, whose daily cycle has no amplitude to take a ratio of on a day without
rain.

Three figures accompany the table, each answering a question the others
average away. The first is the marginal view: every hour on which both sources
report, one quantity to a panel, against the 1:1 line. A constant offset appears
as a cloud displaced parallel to the diagonal, and a source that compresses the
range of a quantity appears as a cloud flatter than the diagonal --- which is
what an amplitude ratio below one looks like before it is reduced to a number.

\begin{figure}[H]
\centering
\includegraphics[width=\textwidth]{PF_F16_cross_source_scatter.png}
\caption{ERA5 against the Meteosystem station, every paired hour, one panel per
shared quantity. The dashed line is 1:1 and the bias and correlation printed in
each panel are the whole-record values of Table~\ref{tab:cross_source}, not a
second computation. Darker hexagons hold more samples, on a logarithmic count
scale. Precipitation is bounded at 20\,mm in the hour, excluding the station's
single heaviest observation.}
\label{fig:cross_source_scatter}
\end{figure}

The second asks whether the disagreement holds still. This is the question that
decides whether a source may substitute for another: a constant offset can be
subtracted and the substitution documented, while an offset that changes with
the season cannot, and the two are indistinguishable in any statistic pooled over
the whole record.

\begin{figure}[H]
\centering
\includegraphics[width=\textwidth]{PF_F17_cross_source_stability.png}
\caption{The difference between the two sources, month by month, one panel per
shared quantity. The band is the interquartile range of the difference within
each month, so a drift and a widening can be told apart. Positive is ERA5
reading high.}
\label{fig:cross_source_stability}
\end{figure}

The third asks about the day. Two sources may agree on the level of a quantity
and still disagree about when it peaks and how far it travels between its
extremes, and neither of the first two figures can show that, because both
average the time of day away.

\begin{figure}[H]
\centering
\includegraphics[width=\textwidth]{PF_F18_cross_source_diurnal.png}
\caption{Mean daily cycle of each shared quantity in each season, both sources
on one axes. The sources are distinguished by line style, since colour carries
the identity of the quantity throughout this report. The two panels of a row
share a vertical scale.}
\label{fig:cross_source_diurnal}
\end{figure}

\begin{table}[H]
\centering
\small
\begin{tabular}{llrrrrrrr}
\toprule
\textbf{Quantity} & \textbf{Split} & \textbf{N} & \textbf{Bias} & \textbf{MAE} &
\textbf{RMSE} & \textbf{r} & \textbf{Amp Ratio} & \textbf{Phase Lag} \\
\midrule
\input{../outputs/PF_T14_cross_source.tex}\\
\bottomrule
\end{tabular}
\caption{ERA5 against the Meteosystem station, over every quantity both sources
report, on their shared hourly record. The station is the reference, so a
positive bias is the reanalysis reading high. Each quantity appears once over the
whole record and once for each season. Units are those of the quantity itself,
given in Tables~\ref{tab:era5_variables} and~\ref{tab:ground_variables}.}
\label{tab:cross_source}
\end{table}

% ===========================================================================
\section{Comparison with the on-structure measurements}
\label{sec:comparison}

The comparison brings together ERA5, the Meteosystem station, and the final
values produced by Study~01. It does not repeat the on-structure data-quality
study. Air temperature is read from its final joined and corrected
\texttt{tair} column. Solar radiation is read from \texttt{n\_sr\_ok}, with all
timestamps on days marked by \texttt{sr\_suspect} excluded. Where Study~01
classified a radiation value as \texttt{night}, the corrected value of
$0\,\mathrm{W/m^2}$ in \texttt{n\_sr\_ok} is used; the original non-zero sensor
floor is not restored.

All three sources will be expressed in common units and placed on an hourly UTC
grid. The twenty-minute on-structure values and the half-hourly Meteosystem
values will be aggregated by hourly mean, while ERA5 is already hourly. Every
valid on-structure value will remain eligible for comparison. Pairwise
statistics will use the overlap with each external source, and direct
three-source figures will use the timestamps on which all three are available.
No proxy value will be interpolated to manufacture an overlap.

\subsection{Air temperature}

The temperature comparison will align the three series over every valid
on-structure observation. It will show their absolute levels, common temporal
variation, and differences. Agreement will be summarised with bias, error,
correlation, and differences in daily amplitude and phase, reported by season
where the coverage supports that division.

\begin{figure}[H]
\centering
\includegraphics[width=\textwidth]{PF_F08_tair_series.png}
\caption{Aligned three-source air-temperature time series over the valid on-structure timestamps.}
\label{fig:tair_series}
\end{figure}

\begin{figure}[H]
\centering
\includegraphics[width=\textwidth]{PF_F09_tair_scatter.png}
\caption{Direct comparison scatter plots of each air-temperature proxy against the on-structure values.}
\label{fig:tair_scatter}
\end{figure}

\begin{figure}[H]
\centering
\includegraphics[width=\textwidth]{PF_F10_tair_diurnal.png}
\caption{Air-temperature diurnal cycles across sources.}
\label{fig:tair_diurnal}
\end{figure}

\begin{figure}[H]
\centering
\includegraphics[width=\textwidth]{PF_F11_tair_stability.png}
\caption{Air-temperature stability and difference analysis.}
\label{fig:tair_stability}
\end{figure}

\begin{table}[H]
\centering
\small
\begin{tabular}{llrrrrrrrrr}
\toprule
\textbf{Source} & \textbf{Season} & \textbf{N} & \textbf{Bias} & \textbf{MAE} & \textbf{RMSE} & \textbf{r} & \textbf{Amp (On)} & \textbf{Amp (Prx)} & \textbf{Amp Ratio} & \textbf{Phase Diff} \\
\midrule
\input{../outputs/PF_T10_tair_agreement.tex}\\
\bottomrule
\end{tabular}
\caption{Air-temperature agreement metrics.}
\label{tab:tair_agreement}
\end{table}

\subsection{Solar radiation}

The radiation comparison will apply the Study~01 validity decisions before any
alignment. It will therefore include corrected night-time zeros, retain valid
daylight measurements, and exclude rejected samples and suspect days. Because
radiation is strongly controlled by time of day and season, agreement will be
examined both over the full valid set and within the seasonal daily cycle.

\begin{figure}[H]
\centering
\includegraphics[width=\textwidth]{PF_F12_sr_series.png}
\caption{Aligned three-source solar-radiation time series over the valid on-structure timestamps.}
\label{fig:sr_series}
\end{figure}

\begin{figure}[H]
\centering
\includegraphics[width=\textwidth]{PF_F13_sr_scatter.png}
\caption{Direct comparison scatter plots of each solar-radiation proxy against the on-structure values.}
\label{fig:sr_scatter}
\end{figure}

\begin{figure}[H]
\centering
\includegraphics[width=\textwidth]{PF_F14_sr_diurnal.png}
\caption{Solar-radiation diurnal cycles across sources.}
\label{fig:sr_diurnal}
\end{figure}

\begin{figure}[H]
\centering
\includegraphics[width=\textwidth]{PF_F15_sr_stability.png}
\caption{Solar-radiation stability and difference analysis.}
\label{fig:sr_stability}
\end{figure}

\begin{table}[H]
\centering
\small
\begin{tabular}{llrrrrrrrrr}
\toprule
\textbf{Source} & \textbf{Season} & \textbf{N} & \textbf{Bias} & \textbf{MAE} & \textbf{RMSE} & \textbf{r} & \textbf{Amp (On)} & \textbf{Amp (Prx)} & \textbf{Amp Ratio} & \textbf{Phase Diff} \\
\midrule
\input{../outputs/PF_T11_sr_agreement.tex}\\
\bottomrule
\end{tabular}
\caption{Solar-radiation agreement metrics.}
\label{tab:sr_agreement}
\end{table}

\subsection{Compatibility and substitution}

The final assessment will state separately whether ERA5 or Meteosystem is a
defensible substitute for on-structure air temperature and for on-structure
solar radiation. A source may be useful without being interchangeable with the
local measurement: a stable bias or amplitude difference may permit a documented
transformation, whereas an unstable seasonal relation, inconsistent phase, or
poor correspondence during the valid local window would rule out substitution.
The verdict will specify the accepted source, any required transformation, the
period over which the evidence applies, and the limitations carried into the
later inclination analysis.

% ===========================================================================
\section{Conclusions}
\label{sec:conclusions}

This section will summarise the character of the two external records and give
the source verdict for each of the two forcings. It will distinguish evidence
that supports filling the missing on-structure coverage from evidence that only
supports using a source as a broader environmental covariate. Any unresolved
metadata or unsupported substitution will remain explicit rather than being
converted into an assumption.

\end{document}
